\Askhsh{33631}{4}
\begin{ylh}{1}
    \item 2.2 Παρουσίαση Στατιστικών Δεδομένων
    \item 2.3 Μέτρα Θέσης και Διασποράς
\end{ylh}
\begin{wrapfigure}[18]{r}{7cm}
\begin{mytblr}{}
Βαθμός $x_i$ & Συχνότητα $\nu_i$ & $x_i \cdot \nu_i$ & $N_i$ \\
6 & 3 & & \\
8 & 2 & & \\
10 & 4 & & \\
11 & 1 & & \\
14 & 5 & & \\
15 & 1 & & \\
18 & 3 & & \\
20 & 1 & & \\
ΣΥΝΟΛΟ & 20 & & \SetCell{bg=gray!30} \\
\end{mytblr}
\end{wrapfigure}
Ένας καθηγητής Μαθηματικών που διδάσκει στο ΕΠΑΛ μιας πόλης, κατέγραψε τη βαθμολογία που πέτυχαν οι 20 μαθητές του, στις Πανελλαδικές Εξετάσεις, στην κλίμακα $0 - 20$. Στη συνέχεια κατασκεύασε το διπλανό πίνακα, όπου σημείωσε τους βαθμούς που πέτυχαν οι μαθητές του (η στήλη των $x_i$) και το πλήθος-συχνότητα των μαθητών που τους πέτυχε (η στήλη των $\nu_i$). $N_i$ είναι η αθροιστική συχνότητα.

\begin{alist}
    \item Να συμπληρώσετε τις δύο στήλες του πίνακα που δεν έχουν συμπληρωθεί. \monades{6}
    \item Να βρείτε τη μέση τιμή της βαθμολογίας, που πέτυχαν οι μαθητές του ΕΠΑΛ αυτού. \monades{7}
    \item Όταν ο σύλλογος των διδασκόντων του ΕΠΑΛ αυτού, συζήτησε τα αποτελέσματα που πέτυχαν οι μαθητές τους, ο καθηγητής των Μαθηματικών υποστήριξε ότι:
    \begin{rlist}
        \item το πανελλαδικό ποσοστό, όσων μαθητών έγραψε κάτω από τη βάση (κάτω από 10) ξεπερνάει το 52\% και είναι υπερδιπλάσιο από το αντίστοιχο ποσοστό των μαθητών του σχολείου.
        \item το πανελλαδικό ποσοστό, όσων μαθητών έγραψε άριστα (από 18 έως 20) είναι μόλις το 5\%, ενώ το αντίστοιχο ποσοστό των μαθητών του σχολείου είναι τετραπλάσιο αυτού.
    \end{rlist}
    Συμφωνείτε με τις απόψεις που διατύπωσε ο Μαθηματικός; Να αιτιολογήσετε την απάντησή σας.
    \monades{12}
\end{alist}

